\chapter{Logger API}

\section{Opening the logger}

\begin{verbatim}
  bool Open(Facility facility = Facility::user,
            const std::vector<MessageSink *> & sinks = {},
            const char *ident = nullptr)

  bool Open(const std::string & facility,
            const std::vector<MessageSink *> & sinks = {},
            const char *ident = nullptr)
\end{verbatim}

If \texttt{sinks} is empty, the logger will send all log messages
to the loopback address where it is expected to be processed by
\textit{mclogd}.

If \texttt{ident} is \texttt{nullptr}, the logger will use the ident
given by the runtime.  On macOS and FreeBSD, this comes from
\texttt{getprogname(3)}.  On Linux, this comes from
\texttt{program\_invocation\_short\_name(3)}.

\section{Alternate or additional sinks}
The logger sends messages to \textit{sinks}, which must implement
the \texttt{Dwm::Mclog::MessageSink} interface.

\subsection{Dwm::Mclog::MessageSink interface}
\begin{verbatim}
  virtual bool Process(const Message & msg)
\end{verbatim}

\subsection{Adding or removing sinks}
Sinks may be added to the logger by using the \texttt{AddSinks()}
member of the logger.

\begin{verbatim}
  bool AddSinks(const std::vector<MessageSink *> & sinks)
\end{verbatim}

Sinks may be removed from the logger by using the \texttt{RemoveSinks()}
member of the logger.

\begin{verbatim}
  bool RemoveSinks(const std::vector<MessageSink *> & sinks)
\end{verbatim}

\subsection{Dwm::Mclog::LogFile}

\texttt{Dwm::Mclog::LogFile} provides a sink that writes logged messages
to a file, with configurable rollover period, permissions and number of
archives to keep.  If all you need is a log file, and/or do not want or
need networked logging, it is a good choice of sink in the list of sinks
given as the second argument of the \texttt{Open()} member of the logger.
It is declared in the \texttt{DwmMclogLogFile.hh} header file.

\begin{verbatim}
  LogFile(const std::string & path, mode_t permissions = 0644,
          RollPeriod period = RollPeriod::days_1, uint32_t keep = 7,
          FileFormat format = FileFormat::text)
\end{verbatim}

\subsection{Dwm::Mclog::OstreamSink}

\begin{verbatim}
  OstreamSink(std::ostream & os)
\end{verbatim}

\texttt{OstreamSink} may be used to create a \texttt{MessageSink} from
a \texttt{std::ostream}.  Note that \texttt{OstreamSink} does not
participate in lifetime management of \texttt{os}.  The caller must
ensure that \texttt{os} outlives the \texttt{OstreamSink}.  It is declared
in the \texttt{DwmMclogOstreamSink.hh} header file.

\subsection{Dwm::Mclog::SyslogSink}

\begin{verbatim}
  SyslogSink(const char *ident, int facility)
\end{verbatim}

\section{Sending log messages}
The \texttt{MCLOG} macro is used to log a message.

\section{Closing the logger}
\begin{verbatim}
  bool Close()
\end{verbatim}

\section{Examples}

\begin{minipage}{\linewidth}
\subsection{Logging to \textit{mclogd}}
The program in listing~\ref{lst:loopback} sends a log message to
\textit{mclogd} on the local host.  This is far and away my most
common usage.

\lstinputlisting[language=c++,caption={\texttt{loopback.cc}},captionpos=b,label={lst:loopback}]{examples/Example01_loopback.cc}
\end{minipage}

\begin{minipage}{\linewidth}
\subsection{Logging to a \texttt{LogFile}}
The program in listing~\ref{lst:logfile} writes a log message to a \texttt{LogFile} in the current directory.

\lstinputlisting[language=c++,caption={\texttt{logfile.cc}},captionpos=b,label={lst:logfile}]{examples/Example02_LogFile.cc}
\end{minipage}

\begin{minipage}{\linewidth}
The program in listing~\ref{lst:cerr} writes a log message to \texttt{std::cerr}.
\subsection{Logging to \texttt{cerr}}
\lstinputlisting[language=c++,caption={\texttt{cerr.cc}},captionpos=b,label={lst:cerr}]{examples/Example03_cerr.cc}
\end{minipage}
